\subsection{Robust root finding with bounds}\label{\refsec:RobustRoot}
In the following subsections, we identify policies that maximize the political objective function by solving the first order condition as a root-finding problem instead of using the political objective function directly. This is paramount to avoid having the entire distribution of past savings as state variables for the political problem. Also, we will generally only work with $\tau \in[0,1]$ for numerical stability. To deal with these two points (that we work with first order conditions and only want to evaluate most functions on $\tau\in[0,1]$), we set up an auxiliary root-finding problem.

Notationally, let $z(\tau) = f\left([\tau]_{[0,1]}\right)$ denote the first order condition evaluated in the interior $[\tau]_{[0,1]} \equiv \min(1, \max(0,\tau))$. Instead of trying to solve for $f(\tau) = 0$, we suggest solving for
\begin{align}\label{\refeq:root}
    h(z,\tau) = z - |z| \underbrace{\left[a_0\min\left(\tau,0\right) + a_1\left(\max(\tau,1)-1\right)\right]}_{\equiv g(\tau)} = 0,
\end{align}
where $a_0, a_1>0$ are some technical parameters. In this case we have three regions:
\begin{align*}
    h\left(z,\tau\right) = \left\lbrace \begin{array}{ll} f(0)-|f(0)| a_0 \tau, & \tau<0 \\
    f(\tau), & \tau\in[0,1] \\
    f(1)-|f(1)|a_1(\tau-1), & \tau>1\end{array}\right., && \dfrac{\partial h}{\partial \tau} = \left\lbrace \begin{array}{ll} -|f(0)|a_0, & \tau<0 \\
    \dfrac{\partial f}{\partial \tau}, & \tau\in[0,1] \\
    -|f(1)|a_1, & \tau>1\end{array}\right.
\end{align*}

If the function $f(\tau)$ is well behaved (i.e.\ identifies a maximum), we can solve simply using $h(z(\tilde{\tau}), \tilde{\tau}) = 0$ and then report the bounded solution as $\tau = [\tilde{\tau}]_{[0,1]}$. As long as the function is well behaved, we can use gradient-based solvers to identify roots.

If the function is not well behaved, we may still use $h(z, \tau)$ to detect boundary solutions in some instances. Assume for instance that $f(\tau) >0$ for all $\tau\in[0,1]$, but that $\partial f/\partial \tau <0$. In this case a gradient-based algorithm would have a hard time identifying a solution starting in the interior, as it would suggest trying out lower values of $\tau$. However, note that for $\tau\geq 1$, the method still works and identifies $\tilde{\tau} = 1+1/a_1$ and thus the boundary solution of $\tau = 1$. Similarly, assume that $f(\tau)<0$ for all $\tau\in[0,1]$ and that $\partial f/\partial \tau >0$. In this case, the same problem exists for a gradient-based solver in the interior, but searching outside the bounds still ensures that we identify the correct root $\tilde{\tau} = -1/a_0$ and thus $\tau = 0$.
